Using LLMs as fixed software components exposes a linguistic attack surface
whose root cause is the native inability to distinguish malicious from
non-malicious inputs at the model's input, from which \textit{prompt injection}
follows. The defense, consolidated both scientifically and commercially in
industrial deployments, places an input \textit{guardrail} before the main
execution. That filter, however, is itself an LLM that receives instruction and
data through the same channel, and inherits the very vulnerability it exists to
mitigate. This work measures what is gained by applying to the filter itself
the separation between instruction channel and data channel that the literature
proposes for the main model, and by having the filter preserve and mark the
suspicious span instead of removing it, forwarding the annotation to the
component that will carry out the task. In essence, the goal is to improve the
security of LLM-based systems without resorting to a \textit{fine-tuned} model,
but by adding security techniques that modify the \textit{input}. The
experimental design is $2 \times 2$: XML delimiting at the filter's input and
selective marking at its output, totaling four conditions, plus a no-filter
test as the initial reference and a \textit{fine-tuned} filter from the same
family, evaluated under the same four conditions, as a point of comparison. The
evaluation uses two open models
without \textit{fine-tuning}, each one instantiating the whole system, over
three public \textit{benchmarks}: BIPIA, SEP and NotInject. There are 42 runs
per repetition, repeated ten times: 420 runs and 131,460 records. The outcome
of each item is its mean over the ten repetitions, measured at the output of
the main component, in attack success rate, utility on benign tasks and
over-blocking, with paired \textit{bootstrap} confidence intervals over the
same items. Combining both factors reduced the attack success rate on
BIPIA by 0.084 for Llama-3.1-8B and by 0.139 for gpt-oss-20b, both with a 95\%
interval away from zero. On the \textit{fine-tuned} filter the same
combination reduced it by 0.139, showing that the gain does not depend on the
filter being a general-purpose model. Marking rather than removing proved a
viable alternative to removal-based defenses, with greater protection than the
\textit{fine-tuned} model from the same family and far easier implementation.
